% FILE: appendix/OC_1_3_3_IDENTITY_RESIDUE_REBIRTH_CATEGORY.tex
% OC Core 1.3.3 scientific closure artifact.

\section{OC 1.3.3 Identity, Residue, and Rebirth Category}
\label{app:oc133-identity-residue-rebirth}

\paragraph{Objects and morphisms.}
Objects are time-sliced carriers, live realizations, and residue structures.
Morphisms include lawful evolution, collapse, residue extraction, embedding
transfer, fusion, split, and rebirth witness maps.

\paragraph{Identity preservation.}
An identity-preserving continuation is a morphism preserving the typed carrier,
required axes, thresholds, cycle support, and admissible live realization.
Death is loss of identity-preserving live morphisms beyond a collapse time.
Rebirth is a later live object connected by a residue-bearing morphism but not
by identity-preserving continuation.

\paragraph{Theorem.}
Continuation, reconstruction, clone, residue, and rebirth are distinct
morphism classes.

\paragraph{Proof.}
The classes differ by preservation predicates.  Continuation preserves live
identity support; reconstruction preserves design information without live
lineage; clone preserves a pattern but not the original carrier; residue lacks
live admissibility; rebirth has residue linkage without identity-preserving
continuation.  Because the predicates are mutually separating, the categories
cannot collapse into narrative equivalence.
